\documentclass[11pt, UTF8]{ctexart}
\usepackage{amsmath, amsfonts, amssymb}
\usepackage{extarrows} 
\usepackage{geometry}                   
\usepackage{xcolor}                     
\usepackage{enumitem}                   
\usepackage[most]{tcolorbox}            
\usepackage{fancyhdr}                   
\setlength{\headheight}{13.6pt}
\geometry{a4paper, margin=0.8in}        
\pagestyle{fancy}
\fancyhf{}
\lhead{\color{fblue} \hwdate 作业 答案} 
\rhead{\color{fblue} \today}
\cfoot{\thepage}

\definecolor{fblue}{RGB}{0, 74, 142}
\definecolor{fred}{RGB}{204, 26, 26}

\newtcolorbox{problem}[2]{
    enhanced,
    colback=white,
    colframe=fblue!80,
    arc=2mm,
    boxrule=0.8pt,
    before upper={\textbf{\color{fblue} 问题 #1 (#2)}\hspace{0.5em}},
    before skip=1.5em,
    after skip=1.5em,
    left=3mm, right=3mm, top=2mm, bottom=2mm
}

\newcommand{\hwdate}{4月2日}
\newcommand{\Prob}{P}
\newcommand{\comp}[1]{\overline{#1}\mskip 3mu}
\newcommand{\ans}[1]{\mathbf{\color{fred}#1}}

\begin{document}

\begin{center}
    {\LARGE \bfseries \color{fblue} \hwdate 作业 答案}
\end{center}

\begin{problem}{1}{2.5}
一房间有 $3$ 扇同样大小的窗子，其中只有一扇是打开的. 有一只鸟自开着的窗子飞入了房间，它只能从开着的窗子飞出去. 鸟在房间里飞来飞去，试图飞出房间. 假定鸟是没有记忆的，它飞向各扇窗子是随机的.
\begin{enumerate}
    \item[(1)] 以 $X$ 表示鸟为了飞出房间试飞的次数，求 $X$ 的分布律.
    \item[(2)] 户主声称，他养的一只鸟是有记忆的，它飞向任一窗子的尝试不多于一次. 以 $Y$ 表示这只聪明的鸟为了飞出房间试飞的次数. 如户主所说是确实的，试求 $Y$ 的分布律.
    \item[(3)] 求试飞次数 $X$ 小于 $Y$ 的概率和试飞次数 $Y$ 小于 $X$ 的概率.
\end{enumerate}
\end{problem}

\paragraph{\color{fblue}解} 
（1）本题的试飞次数是指记录鸟儿飞向窗子的次数加上最后飞离房间的一次，其分布律为
\[ P\{X = k\} = \ans{\left(\frac{2}{3}\right)^{k-1} \frac{1}{3}}, \quad k = 1, 2, \cdots \]

（2）由题意 $Y$ 的可能值为 $1, 2, 3$. \\
$\{Y=1\}$ 表明鸟儿从 3 扇窗子中选对了一扇，因对鸟儿而言，3 扇窗是等可能被选取的，故 $P\{Y=1\}=\frac{1}{3}$. \\
$\{Y=2\}$ 表明第一次试飞失败（选错了窗子），失败方式有 2，故第一次失败的概率为 $\frac{2}{3}$；第二次，鸟儿舍弃已飞过的那扇窗，而从剩下的一开一关的两窗选一，成功机会为 $\frac{1}{2}$，故 $P\{Y=2\} = \frac{2}{3} \times \frac{1}{2} = \frac{1}{3}$. \\
对有记忆鸟儿来说，$\sum_{i=1}^{3} P\{Y=i\} = 1$，故 $P\{Y=3\} = \frac{1}{3}$. \\
因此 $Y$ 的分布律为
\[ P\{Y = i\} = \ans{\frac{1}{3}}, \quad i = 1, 2, 3. \]

（3）(i) $\{X<Y\}$ 可分解为下列 3 个两两互不相容的事件之和，即
\[ \{X < Y\} = \{(X = 1) \cap (Y = 2)\} \cup \{(X = 1) \cap (Y = 3)\} \cup \{(X = 2) \cap (Y = 3)\}, \]
故
\begin{align*}
    P\{X < Y\} &= P\{(X = 1) \cap (Y = 2)\} + P\{(X = 1) \cap (Y = 3)\} \\
    &\quad + P\{(X = 2) \cap (Y = 3)\}.
\end{align*}
因为两只鸟儿的行动是相互独立的，从而
\begin{align*}
    P\{X < Y\} &= P\{X=1\}P\{Y=2\} + P\{X=1\}P\{Y=3\} + P\{X=2\}P\{Y=3\} \\
    &= \frac{1}{3} \times \frac{1}{3} + \frac{1}{3} \times \frac{1}{3} + \frac{2}{9} \times \frac{1}{3} = \ans{\frac{8}{27}}.
\end{align*}

(ii)
\begin{align*}
    P\{Y < X\} &= 1 - P\{X < Y\} - P\{X = Y\} \\
    &= 1 - \frac{8}{27} - \sum_{k=1}^{3} P\{(X = k) \cap (Y = k)\} \\
    &= 1 - \frac{8}{27} - \sum_{k=1}^{3} P\{X = k\}P\{Y = k\} \\
    &= 1 - \frac{8}{27} - \frac{1}{3} \times \frac{1}{3} - \frac{2}{9} \times \frac{1}{3} - \frac{4}{27} \times \frac{1}{3} = \ans{\frac{38}{81}}.
\end{align*}

\vspace{1.5em}

\begin{problem}{2}{2.14}
某人家中在时间间隔 $t$（以 h 计）内接到电话的次数 $X$ 服从参数为 $2t$ 的泊松分布.
\begin{enumerate}
    \item[(1)] 若他外出计划用时 10 min，问其间电话铃响一次的概率是多少？
    \item[(2)] 若他希望外出时没有电话的概率至少为 0.5，问他外出应控制的最长时间是多少？
\end{enumerate}
\end{problem}

\paragraph{\color{fblue}解}
以 $X$ 表示此人外出时电话铃响的次数，则 $X \sim \pi(2t)$，其中 $t$ 表示外出的总时间，即 $X$ 的分布律为 $P\{X=k\} = (2t)^k \mathrm{e}^{-2t}/k!, k=0,1,2,\cdots$

（1）$t = 10/60 = 1/6$ 时，$X \sim \pi\left(2 \times \frac{1}{6}\right)$，故所求概率为
\[ P\{X=1\} =\ans{ \frac{1}{3} \mathrm{e}^{-1/3} }= \ans{0.238 \,8}. \]

（2）设外出最长时间为 $t(\mathrm{h})$，因 $X \sim \pi(2t)$，无电话打进的概率为
\[ P\{X=0\} = \mathrm{e}^{-2t}, \]
要使 $P\{X=0\} = \mathrm{e}^{-2t} \geqslant 0.5$，即要使 $\mathrm{e}^{2t} \leqslant 2$，由此得
\[ t \leqslant \ans{\frac{1}{2} \ln 2} = \ans{0.346\,6} \,(\mathrm{h}), \]
即外出时间应控制少于 $\ans{20.79} \min$.

\vspace{1.5em}

\begin{problem}{3}{2.15}
保险公司在一天内承保了 5 000 张相同年龄、为期一年的寿险保单，每人一份. 在合同有效期内若投保人死亡，则公司需赔付 3 万元. 设在一年内，该年龄段的死亡率为 0.001 5，且各投保人是否死亡相互独立. 求该公司对于这批投保人的赔付总额不超过 30 万元的概率（利用泊松定理计算）.
\end{problem}

\paragraph{\color{fblue}解}
设这批投保人在一年内死亡的人数为 $X$，则 $X \sim b(5\,000, 0.001\,5)$，因每死亡一人公司需赔付 3 万元，故公司赔付不超过 30 万元意味着在投保期内死亡人数不超过 $30/3=10$（人），从而所求概率为
\[ P\{X \leqslant 10\} = \sum_{k=0}^{10} \binom{5\,000}{k} 0.001\,5^k (1 - 0.001\,5)^{5\,000-k}. \]
若用泊松定理近似可以认为 $X \sim \pi(7.5)$，于是
\[ P\{X \leqslant 10\} \approx \sum_{k=0}^{10} \frac{7.5^k \mathrm{e}^{-7.5}}{k!} \xlongequal{\text{查表}} \ans{0.862\,2}. \]

\end{document}